% ---
% id: protocol-compiler/rem:coin-scope
% kind: rem
% title: Scope of the private-coin assumption
% summary: >-
%   A state-dependent marginal of a setting-conditioned public choice is not a
%   violation of [[A-COIN]]: it is the output-state model of the reductions
%   module. A violation is repaired by no encoding, only by a different protocol
%   description.
% prominence: standard
% status: accepted
% uses: [def:weighted]
% mentions: [ass:coin, lem:weighted, thm:choice-encoding, rem:relabel, reductions/ass:tsm]
% assumes: []
% refs: []
% provenance: newly-derived
% lean:
% ---

\begin{remark}\label{rem:coin-scope}
No encoding repairs a violation of Assumption~\ref{ass:coin}; the remedy is a
different protocol description. Under a violation the weighted POVM
misdescribes the round: for the \tref{setting-independent form}
\eqref{eq:choicemass} no longer gives the physical choice statistics, and for
the setting-dependent form the relabelled family of Remark~\ref{rem:relabel}
no longer describes how $f$ was drawn; a choice correlated with the adversary,
or with the device outcome beyond its setting, is exactly the attack the
encoding assumes away. The state-dependent \emph{marginal} of $f$
that the setting-dependent form of Definition~\ref{def:weighted} produces is
not a violation: it is the output-state model of
\xref{reductions/ass:tsm}, and
Lemma~\ref{lem:weighted}\,(iii) is asserted for the setting-independent form
only (Theorem~\ref{thm:choice-encoding}).
\end{remark}
